% Delta de §10.27. Insertar al final de 05_ciclo_catalan.tex.
% La actualización afín y su balance permanecen en sus propietarios originales.
\subsection{Trazas de retorno y momentos del cociente constitutivo}
\label{amp-afin-sec}

La actualización del registro y la respuesta del ciclo comparten el
transporte dodecafásico y conservan información de distinto tipo.
El evento firmado se obtiene de la dirección y la orientación de cada
ventana de la historia APP--TRIT--TPK. Su incremento vectorial y el
balance~\eqref{mem:balance} retienen la interacción con el registro previo.
La construcción siguiente prolonga la lectura espectral del mismo ciclo,
una vez fijados sus subespacios de incidencia. Este orden permite componer
retorno, determinante y momentos conservando la historia que los precede.

Conservamos el operador $C$ y los proyectores de
\eqref{iii:eq:projectors34}, y definimos
\begin{equation}
 E_{\rm cyc}=P_3^{\mathrm{cyc}}-P_4^{\mathrm{cyc}},\qquad
 b_n=\operatorname{Tr}(C^n E_{\rm cyc}),\qquad n\geq0.
 \label{amp-afin-trazas}
\end{equation}
El símbolo $b_n$ distingue esta traza espectral del evento firmado
$a_m$ de una ventana. La traza utiliza el operador de transporte; el
evento utiliza además la trayectoria efectivamente recorrida.

\begin{proposition}[Tabla de trazas producida por la incidencia]
\label{amp-afin-prop-tabla}
Se cumplen
\begin{equation}
 \begin{gathered}
 b_n=3\mathbf1_{3\mid n}-4\mathbf1_{4\mid n},\\
 (b_1,\ldots,b_{12})=(0,0,3,-4,0,3,0,-4,3,0,0,-1).
 \end{gathered}
 \label{amp-afin-tabla}
\end{equation}
El período mínimo es doce y la suma en un período es cero. Para
$B_N=\sum_{n=1}^{N}b_n$,
\begin{equation}
 B_N=3\lfloor N/3\rfloor-4\lfloor N/4\rfloor,\qquad |B_N|\leq3.
 \label{amp-afin-parciales}
\end{equation}
\end{proposition}
\begin{proof}
En $H_d$, el operador $C$ permuta cíclicamente las $d$ clases de
coordenadas. Su potencia $n$ tiene $d$ puntos fijos si $d$ divide a $n$
y ninguno en caso contrario. Como $P_d^{\mathrm{cyc}}$ conmuta con $C$,
$\operatorname{Tr}(C^nP_d^{\mathrm{cyc}})$ es la traza de esa restricción.
Restar los casos $d=3,4$ prueba la fórmula. El mínimo período positivo
es doce: $b_n=-1$ ocurre exactamente en los múltiplos de doce.
En un período las contribuciones son $3\cdot4-4\cdot3=0$.
La suma finita cuenta los múltiplos y da~\eqref{amp-afin-parciales}.
Sus valores para $N=0,\ldots,11$ son
$(0,0,0,3,-1,-1,2,2,-2,1,1,1)$ y se repiten al sumar doce.
\end{proof}

\begin{theorem}[Reconstrucción logarítmica del lector constitutivo]
\label{amp-afin-thm-logdet}
Para $|s|<1$, el cociente $f$ de~\eqref{iii:eq:trace-det} verifica
\begin{align}
 -\log f(s)&=\sum_{n\geq1}\frac{b_n}{n}s^n,\label{amp-afin-logdet}\\
 -\frac{s f'(s)}{f(s)}
 &=\frac{3s^3}{1-s^3}-\frac{4s^4}{1-s^4}.
 \label{amp-afin-logder}
\end{align}
La rama holomorfa queda fijada por $\log f(0)=0$.
En $0<s<1$ coincide con el logaritmo real.
\end{theorem}
\begin{proof}
Los polinomios $1-s^3$ y $1-s^4$ no se anulan en el disco abierto.
La expansión $-\log(1-w)=\sum_{k\geq1}w^k/k$, obtenida integrando
la serie geométrica desde cero, converge uniformemente en discos
compactos. Con $w=s^3,s^4$, su diferencia da
$-\log f(s)=\sum_{k\geq1}s^{3k}/k-\sum_{k\geq1}s^{4k}/k$.
El coeficiente de $s^n$ es $b_n/n$.
La convergencia local uniforme permite derivar y sumar las dos series
geométricas, lo que prueba~\eqref{amp-afin-logder}.
La condición en cero fija la constante de integración y la rama.
\end{proof}

Así, $b_0/12=-1/12$ registra el defecto normalizado de dimensión,
mientras toda la sucesión $b_n$ reconstruye una función.
Son dos lecturas de distinto alcance del mismo par de subespacios.
Al sustituir $s=q_\pm^{30}$, los argumentos siguen siendo las salidas
de la construcción angular previa; la expansión no los selecciona.

\begin{theorem}[Momentos de profundidad y representación de Mellin]
\label{amp-afin-thm-momentos}
La serie
\begin{equation}
 M_b(z)=\sum_{n\geq1}\frac{b_n}{n^z}
 \label{amp-afin-momento}
\end{equation}
converge localmente de manera uniforme para $\operatorname{Re}z>0$.
Si $\sigma=\operatorname{Re}z>0$, su resto satisface
\begin{equation}
 \left|M_b(z)-\sum_{n=1}^{N}\frac{b_n}{n^z}\right|
 \leq3\left(1+\frac{|z|}{\sigma}\right)(N+1)^{-\sigma}.
 \label{amp-afin-cota}
\end{equation}
En el dominio de convergencia absoluta $\operatorname{Re}z>1$,
\begin{equation}
 M_b(z)=(3^{1-z}-4^{1-z})\zeta(z),\qquad
 \zeta(z)=\sum_{n\geq1}n^{-z}.
 \label{amp-afin-reconocimiento}
\end{equation}
Para $\operatorname{Re}z>0$ se tiene además
\begin{equation}
 \Gamma(z)M_b(z)=
 \int_0^\infty t^{z-1}
 \left(\frac3{e^{3t}-1}-\frac4{e^{4t}-1}\right)dt,
 \label{amp-afin-mellin}
\end{equation}
donde $\Gamma(z)=\int_0^\infty u^{z-1}e^{-u}\,du$.
\end{theorem}
\begin{proof}
La sumación por partes y $|B_n|\leq3$ dan, al pasar al límite superior,
\[
 \sum_{n>N}b_n n^{-z}
 =-B_N(N+1)^{-z}
   +\sum_{n>N}B_n\bigl(n^{-z}-(n+1)^{-z}\bigr).
\]
La diferencia de potencias está acotada en módulo por
$|z|\int_n^{n+1}x^{-\sigma-1}dx$. Esto prueba
\eqref{amp-afin-cota}, la convergencia y su uniformidad en compactos.
Para $\sigma>1$, la convergencia absoluta permite agrupar
los múltiplos de tres y de cuatro por separado, dando
\eqref{amp-afin-reconocimiento}.

La serie geométrica del transporte, evaluada en $s=e^{-t}$, produce
\[
 \sum_{n\geq1}b_n e^{-nt}
 =\frac3{e^{3t}-1}-\frac4{e^{4t}-1}=:K_b(t).
\]
Para $\sigma>1$, se puede integrar término a término porque
$\sum|b_n|n^{-\sigma}<\infty$. La sustitución $u=nt$ da
$\Gamma(z)n^{-z}$ y prueba~\eqref{amp-afin-mellin} en ese semiplano.
En el origen se cancelan los términos $1/t$:
$K_b(t)=1/2+O(t)$. En el infinito, $K_b(t)=O(e^{-3t})$.
La integral es, por ello, holomorfa para $\operatorname{Re}z>0$
por convergencia local uniforme. La serie también lo es por la cota
anterior; el principio de identidad prolonga la igualdad a ese dominio.
\end{proof}

En $0<\operatorname{Re}z\leq1$, el orden creciente de profundidad
fija la lectura convergente de la serie. La realización como traza
absoluta del operador diagonal corresponde a $\operatorname{Re}z>1$.

\begin{corollary}[Evaluaciones compatibles con el extremo constitutivo]
\label{amp-afin-cor-extremos}
Se cumplen
\begin{equation}
 M_b(2)=\frac{\zeta(2)}{12},\qquad
 M_b(1)=\log(4/3)=-\log f(1),\qquad f(1)=\frac34.
 \label{amp-afin-extremos}
\end{equation}
\end{corollary}
\begin{proof}
La primera igualdad resulta de $3^{-1}-4^{-1}=1/12$.
Para la segunda, los parciales exactos son
$H_{\lfloor N/3\rfloor}-H_{\lfloor N/4\rfloor}$, con
$H_j=\sum_{k=1}^j1/k$ y $H_0=0$.
La comparación integral de $1/x$ entre esos dos índices da el límite
$\log(4/3)$. Cancelar $1-s$ antes de evaluar el cociente determinantal
da $f(1)=3/4$.
\end{proof}

El registro afín retiene los eventos y su interacción con la memoria
previa; la tabla de trazas retiene la incidencia de los dos subespacios
a todas las potencias del transporte; el momento aplica una ponderación
de profundidad a esa tabla ya producida. Su periodicidad espectral
convive con una historia de eventos que continúa después del retorno
de fase. Invertir $C$ conserva $b_n$, porque $d\mid n$ es invariante
bajo $n\mapsto-n$; el coeficiente orientado del resolvente de
\eqref{iii:eq:resolvent}, con sus marcas fijas, cambia de signo.
La lectura conjunta conserva una distinción que el determinante
escalar por sí solo no registra. La forma cuadrática y la suma de
coordenadas del registro mantienen sus tipos matemáticos; esta
composición no les asigna unidades de energía ni de carga eléctrica.

\subsection{Composición residual y producto de momentos}
\label{amp-afin-sec-productos}

La composición posterior de lectores conserva también la operación
efectuada sobre la profundidad. En la realización
$\mathcal H=\ell^2(\mathbb N_{>0})$, fijamos
$N\delta_n=n\delta_n$. Sean $a,c$ dos tablas residuales periódicas
ya producidas y $Q_a\delta_n=a(n)\delta_n$,
$Q_c\delta_n=c(n)\delta_n$. Al prolongar ambas al período común,
los dos lectores actúan sobre el mismo índice. Para
$\operatorname{Re}z>1$, escribimos
$M_a(z)=\operatorname{Tr}(N^{-z}Q_a)$ y análogamente $M_c$.

\begin{proposition}[Composición diagonal y composición tensorial]
\label{amp-afin-prop-productos}
Se tienen las identidades
\begin{align}
 \operatorname{Tr}(N^{-z}Q_aQ_c)
 &=\sum_{n\geq1}\frac{a(n)c(n)}{n^z},
 \label{amp-afin-diagonal}\\
 M_a(z)M_c(z)
 &=\operatorname{Tr}_{\mathcal H\otimes\mathcal H}
 \bigl((N\otimes N)^{-z}(Q_a\otimes Q_c)\bigr)\notag\\
 &=\sum_{k\geq1}\frac{(a*c)(k)}{k^z},
 \label{amp-afin-tensor}\\
 (a*c)(k)&=\sum_{d\mid k}a(d)c(k/d),
 \qquad \operatorname{Re}z>1.
 \label{amp-afin-convolucion}
\end{align}
Todos los operadores que figuran dentro de estas trazas son de clase traza.
\end{proposition}
\begin{proof}
El producto diagonal actúa sobre $\delta_n$ mediante
$a(n)c(n)$, lo que prueba~\eqref{amp-afin-diagonal}.
En $\delta_m\otimes\delta_n$, el operador $N\otimes N$
tiene valor propio $mn$ y $Q_a\otimes Q_c$ actúa mediante
$a(m)c(n)$. Si $\sigma=\operatorname{Re}z>1$, la suma de los
módulos de estos coeficientes espectrales está acotada por
$\|a\|_\infty\|c\|_\infty(\sum n^{-\sigma})^2<\infty$.
Esta cota prueba la condición de clase traza y permite factorizar
la suma doble o agruparla por $k=mn$. Los pares con producto $k$
son exactamente $(d,k/d)$, con $d\mid k$.
La cota análoga con una sola suma prueba el caso diagonal.
\end{proof}

En los lectores de esta sección, tomar $a(n)=b_n$ y
$c(n)=\langle v,C^nu\rangle$ da, para $n=3$,
$a(3)c(3)=-3$, mientras $(a*c)(3)=3$.
Esta diferencia de coeficientes exhibe la información que conserva
cada operación. La composición diagonal mantiene una profundidad común;
la tensorial conserva dos profundidades antes de agrupar por su producto.
La realización tensorial no atribuye independencia estadística a las
historias HMT que produjeron los lectores. Las derivadas regularizadas y
los cambios de normalización de otros momentos constituyen operaciones
posteriores distintas de las dos composiciones aquí demostradas.
